<<<<<<< HEAD
% This beamer template was adapted from "beamerthemeUppsala version 2008/05/15"
% for GSHS students, by Chinook Mok and GSHS TeX Society(http://latex.gs.hs.kr).
\documentclass[11pt]{beamer}
\hypersetup{pdfencoding=auto}

\usetheme{Berkeley}

\usepackage{xcolor}
\usepackage{amssymb,amsmath,graphicx}
\usepackage[hangul]{kotex}
\usepackage{multimedia}
\usepackage{setspace}
\usepackage{multicol}
%\usepackage{verbatim} %use 
\usepackage{lipsum}
\graphicspath{{images/}}
%% Footnote Setting %%

\let\oldfootnote\footnote
\renewcommand\footnote[1][]{\oldfootnote[frame,#1]}

\makeatletter % Makes maximum footnote counter of fnsymbol from 9 to 12.
\def\@fnsymbol#1{\ensuremath{\ifcase#1\or *\or \dagger\or \ddagger\or
		\mathsection\or \mathparagraph\or \|\or **\or \dagger\dagger
		\or \ddagger\ddagger\or \mathsection\mathsection
		\or \mathparagraph\mathparagraph \or \|\|\else\@ctrerr\fi}}
\makeatother
\renewcommand{\thefootnote}{\fnsymbol{footnote}}

%% Color Setting %%
\definecolor{Grey}{rgb}{0.6,0.6,0.6}
\definecolor{GSHSred}{RGB}{105,0,0}
\definecolor{GSHSblue}{RGB}{0,0,105}
\definecolor{GSHSRED}{RGB}{80,0,0}
\definecolor{GSHSBLUE}{RGB}{0,0,80}
\definecolor{gshsred}{RGB}{240,210,210}
\definecolor{gshsblue}{RGB}{210,210,240}
\setbeamercolor{title}{bg=white,fg=GSHSblue}
\setbeamercolor{author}{fg=GSHSblue}
\setbeamercolor{institute}{fg=GSHSblue}
\setbeamercolor{date}{fg=GSHSblue}
\setbeamercolor{logo}{bg=GSHSblue}
\setbeamercolor{sidebar}{bg=GSHSblue}
\setbeamercolor{frametitle}{bg=white,fg=GSHSblue}
\setbeamercolor{section in sidebar}{use=sidebar,bg=white,fg=sidebar.bg}
\setbeamercolor{subsection in sidebar}{parent=section in sidebar}
\setbeamercolor{subsubsection in sidebar}{parent=subsection in sidebar}
\setbeamercolor{item}{fg=GSHSblue}
\setbeamercolor{block title}{bg=GSHSBLUE,fg=white}
\setbeamercolor{block body}{bg=gshsblue,fg=black}

%% Frametitle Setting %%
\setbeamerfont{frametitle}{series=\bfseries}
\setbeamertemplate{frametitle}
{ \vspace*{-10mm}
  \leavevmode
  \hspace*{3pt}
  \begin{beamercolorbox}[wd=\paperwidth,ht=1ex,dp=1ex]{frametitle}
    \hspace*{7pt}\underline{\makebox[0.6\paperwidth][l]{
    \Large{\insertframetitle}}}
  \end{beamercolorbox}
}

             
%% Title Page Setting %%
\setbeamerfont{title}{series=\bfseries}
\setbeamertemplate{title page}{
  %% Background Logo
  \begin{picture}(0,0)%
    \setlength{\unitlength}{1cm}% default
    \protect\put(1.3,0.8){%
    \begin{picture}(6,6)(4,10)%
    %%%%%%%%%%%%%%%%%%%%%%%%%%%%%%%%%%%%%%%%%%%%%
      %\includegraphics[width=0.4\paperwidth]{New_Profile_Picture.pdf}%
      \includegraphics[width=0.7\paperwidth]{koistudy_icon.pdf}%
    \end{picture}%
    }%
  \end{picture}%
  \vfill
  \vspace*{10mm}
  \raggedleft
  %% Title
  \usebeamerfont{title}\usebeamercolor[fg]{title}\inserttitle\par
  \vskip 2mm
  %% Subtitle
  \ifx\insertsubtitle\@empty
  \else\usebeamerfont{subtitle}\usebeamercolor[fg]{subtitle}\insertsubtitle
  \fi

  \vskip 3mm
  %% Horizontal line
  \usebeamercolor[fg]{title}\hrule height 2pt\hfill
  \vskip 10mm
  %% Author
  \usebeamercolor[fg]{author}\usebeamerfont{author}\insertauthor
  \vskip 1cm
  %% Institute
  \usebeamercolor[fg]{institute}\usebeamerfont{institute}\insertinstitute
  \vskip 1cm
  %% Date
  \usebeamercolor[fg]{date}\usebeamerfont{date}\insertdate
  \vfill
}

%% itemize bullet setting %%
\setbeamertemplate{itemize items}[ball]

%% block setting %%
\setbeamertemplate{blocks}[rounded][shadow=true]

%% caption setting %%
\setbeamertemplate{caption}[numbered]

%% Title, Author, Institute, Date %%
\title[O주차 - X, X]{O주차 - X, X}
\subtitle[]{2016 이산구조}
\author[]{gs1XXXX(XX기 XXX)}
\institute[GSHS]{경기과학고등학교}
\date[]{\today}

%%%%%%%%%%%%%%%%%%%%%%%%%%%%%%%%%%%%%%%%%
% insert your own profile picture
%\logo{\includegraphics[width=15mm]{stackoverflow.pdf}}
%\logo{\includegraphics[width=8mm]{Profile_Picture.pdf}}

%% Main %%

\begin{document}

\begin{frame}[plain]
\titlepage
\end{frame}
\setbeamertemplate{sidebar left}[sidebar theme]

\setstretch{1.3} %줄간격


\section{1}
\begin{frame}{1.1}\footnotesize

	관성계 : 절대 좌표계 $ (x, y, z)$ \\
	비관성계 : 회전좌표계 $ (x^{\prime}, y^{\prime}, z^{\prime})$ 
	극좌표 $ (r, \theta, z)$\\
	\\
	$ (x, y, z) 	\rightarrow (r, \theta, z)$\\
	수평 방향은 정역학 평형 상태에 있으므로, 	
	$ (x, y) 	\rightarrow (r, \theta)$\\
	
	$ \overrightarrow {F} = F_{x} \hat{i}  + F_{y} \hat{j} $\\
	
	$ F_{x} = m \dfrac{d^{2}x}{dt^{2}}$
	$ F_{y} = m \dfrac{d^{2}y}{dt^{2}}$\\
	
	$ (x, y) = (r \cos \theta, r \sin \theta)$\\
	
	$ F_{r} = F_{x} \cos \theta + F_{y} \sin \theta 
	= m \left ( \dfrac{d^{2}x}{dt^{2}} \cos \theta + \dfrac{d^{2}y}{dt^{2}} \sin \theta \right) $\\
	
	$ F_{\theta} = F_{y} \cos \theta - F_{x} \cos \theta 
	= m \left ( \dfrac{d^{2}y}{dt^{2}} \cos \theta - \dfrac{d^{2}x}{dt^{2}} \cos \theta \right) $\\
	
\end{frame}

\begin{frame}{1.2}\footnotesize

	$ x = r \cos \theta $ \\
	
	$\dfrac{dx}{dt} = \cos \theta \dfrac{dr}{dt} - r \sin \theta \dfrac{d\theta}{dt}$ \\
	
	$\dfrac{d^{2}x}{dt^{2}} = \cos \theta \dfrac{d^{2}r}{dt^{2}} - \sin \theta \dfrac{dr}{dt} - \sin \theta \dfrac{dr}{dt} \dfrac{d\theta}{dt} -r \cos \theta \dfrac{d^{2}\theta}{dt^{2}}$\\
	\\
	$ y = r \sin \theta $ \\
	
	$\dfrac{dy}{dt} = \sin \theta \dfrac{dr}{dt} + r \cos \theta \dfrac{d\theta}{dt}$ \\

	$\dfrac{d^{2}t}{dt^{2}} = \sin \theta \dfrac{d^{2}r}{dt^{2}} + \cos \theta \dfrac{dr}{dt} + \cos \theta \dfrac{dr}{dt} \dfrac{d\theta}{dt} -r \sin \theta \dfrac{d^{2}\theta}{dt^{2}}$\\
	\\

\end{frame}

\begin{frame}{1.3}\footnotesize

	$ F_{r} = F_{x} \cos \theta + F_{y} \sin \theta 
	= m \left ( \dfrac{d^{2}x}{dt^{2}} \cos \theta + \dfrac{d^{2}y}{dt^{2}} \sin \theta \right) $\\
	
	$ F_{\theta} = F_{y} \cos \theta - F_{x} \cos \theta 
	= m \left ( \dfrac{d^{2}y}{dt^{2}} \cos \theta - \dfrac{d^{2}x}{dt^{2}} \cos \theta \right) $\\
	
	$\dfrac{d^{2}x}{dt^{2}} = \cos \theta \dfrac{d^{2}r}{dt^{2}} - \sin \theta \dfrac{dr}{dt} - \sin \theta \dfrac{dr}{dt} \dfrac{d\theta}{dt} -r \cos \theta \dfrac{d^{2}\theta}{dt^{2}}$\\
	
	$\dfrac{d^{2}t}{dt^{2}} = \sin \theta \dfrac{d^{2}r}{dt^{2}} + \cos \theta \dfrac{dr}{dt} + \cos \theta \dfrac{dr}{dt} \dfrac{d\theta}{dt} -r \sin \theta \dfrac{d^{2}\theta}{dt^{2}}$\\
	
	
	정리하면, \\
	
	$ F_{r} = m \left[ \dfrac{d^{2}r}{dt^{2}} - r \left( {\dfrac{d \theta}{dt}} \right)^{2} \right] $\\
	\\
	$ F_{\theta} = m \left[ r \dfrac{d^{2}\theta}{dt^{2}} + 2 \dfrac{dr}{dt} \dfrac{d\theta}{dt}  \right] $\\

\end{frame}


\begin{frame}{1.4}\footnotesize

	$\dfrac{d \theta}{dt} = \Omega $, $ r\dfrac{d \theta}{dt} = r \Omega $, $u_{\theta} $, $\dfrac{dr}{dt} = v_{r}$\\
	
	$\dfrac{d}{dt} r^{2}\Omega = 2 r \dfrac{dr}{dt} \Omega + r^{2} \dfrac{d\Omega}{dt}\\
	=  r \left( r \dfrac{d\Omega}{dt} +2 \dfrac{dr}{dt} \Omega \right) \\
	\\
	$ r F_{\theta} = m \dfrac{d}{dt} \left(r^{2} \Omega \right) $
	$ r^{2} \Omega = const$\\
	\\
	
	$ r \left( r \Omega \right) = r u_{\theta} = const $ \\
	\\
	
	$ R_{1} v_{1} = R_{2} v_{2}$\\
	\\


\end{frame}



\begin{frame}{1.5}\footnotesize

$ dV = dx \cdot dy \cdot dz $\\
$x$ 방향 \\
$ F_{x} = P \cdot \Delta y \cdot \Delta z - \left( P + \Delta P \right) \Delta y \cdot \Delta z$\\
$ F_{x} = - \Delta P \cdot \Delta y \cdot \Delta z $\\

$ \dfrac { \Delta y}{\Delta x } = \dfrac {f\left(x + \Delta x \right) - f\left(x \right)}{ \Delta x}$\\

$f^{\prime} \left(x \right) = lim \dfrac { \Delta y}{\Delta x } = \dfrac {f\left(x + \Delta x \right) - f\left(x \right)}{ \Delta x}$\\

$z = f \left( x, y \right) $, $ y = b \left( 고정 \right) $ //

$\dfrac{\partial z}{\partial x} = lim \dfrac { \Delta z}{\Delta x } = \dfrac {f\left(x + \Delta x, b \right) - f\left(x, b \right)}{ \Delta x}$\\

$\dfrac{\partial z}{\partial y} = lim \dfrac { \Delta z}{\Delta y } = \dfrac {f\left(y + \Delta y, b \right) - f\left(y, b \right)}{ \Delta y}$\\

$ \Delta z = \dfrac{\partial z}{\partial x} \Delta x + \dfrac{\partial z}{\partial y} \Delta y $\\
$ dz = \dfrac{\partial z}{\partial x} dx + \dfrac{\partial z}{\partial y} dy $\\
 \end{frame}

\begin{frame}{1.5}\footnotesize

$ \Delta T = \dfrac{\partial T}{\partial t} \Delta t + \dfrac{\partial T}{\partial x} \Delta x + \dfrac{\partial T}{\partial y} \Delta y + \dfrac{\partial T}{\partial z} \Delta z $\\
$ F_{x} = - \Delta P \cdot \Delta y \cdot \Delta z = \dfrac{\partial P}{\partial x} \cdot \Delta x \cdot \Delta y \cdot \Delta z $\\
$ \rho = \dfrac {m}{\Delta x \cdot \Delta y \cdot \Delta z} $ \\

$ \dfrac {F_{x}}{m} = - \dfrac{1}{\rho} \dfrac{\partial P}{\partial x} $ \\

$ \dfrac {F}{m} = - \dfrac{1}{\rho} \left( \dfrac{\partial P}{\partial x} i + \dfrac{\partial P}{\partial y} j + \dfrac{\partial P}{\partial z} k \right) = - \dfrac{1}{\rho} \nabla P$ \\


\end{frame}

\begin{frame}{4}
	
	\begin{block}{asdf}
		Lorem Ipsum Dolor Sit Amet
	\end{block}
\end{frame}
=======
% This beamer template was adapted from "beamerthemeUppsala version 2008/05/15"
% for GSHS students, by Chinook Mok and GSHS TeX Society(http://latex.gs.hs.kr).
\documentclass[11pt]{beamer}
\hypersetup{pdfencoding=auto}

\usetheme{Berkeley}

\usepackage{xcolor}
\usepackage{amssymb,amsmath,graphicx}
\usepackage[hangul]{kotex}
\usepackage{multimedia}
\usepackage{setspace}
\usepackage{multicol}
%\usepackage{verbatim} %use 
\usepackage{lipsum}
\graphicspath{{images/}}
%% Footnote Setting %%

\let\oldfootnote\footnote
\renewcommand\footnote[1][]{\oldfootnote[frame,#1]}

\makeatletter % Makes maximum footnote counter of fnsymbol from 9 to 12.
\def\@fnsymbol#1{\ensuremath{\ifcase#1\or *\or \dagger\or \ddagger\or
		\mathsection\or \mathparagraph\or \|\or **\or \dagger\dagger
		\or \ddagger\ddagger\or \mathsection\mathsection
		\or \mathparagraph\mathparagraph \or \|\|\else\@ctrerr\fi}}
\makeatother
\renewcommand{\thefootnote}{\fnsymbol{footnote}}

%% Color Setting %%
\definecolor{Grey}{rgb}{0.6,0.6,0.6}
\definecolor{GSHSred}{RGB}{105,0,0}
\definecolor{GSHSblue}{RGB}{0,0,105}
\definecolor{GSHSRED}{RGB}{80,0,0}
\definecolor{GSHSBLUE}{RGB}{0,0,80}
\definecolor{gshsred}{RGB}{240,210,210}
\definecolor{gshsblue}{RGB}{210,210,240}
\setbeamercolor{title}{bg=white,fg=GSHSblue}
\setbeamercolor{author}{fg=GSHSblue}
\setbeamercolor{institute}{fg=GSHSblue}
\setbeamercolor{date}{fg=GSHSblue}
\setbeamercolor{logo}{bg=GSHSblue}
\setbeamercolor{sidebar}{bg=GSHSblue}
\setbeamercolor{frametitle}{bg=white,fg=GSHSblue}
\setbeamercolor{section in sidebar}{use=sidebar,bg=white,fg=sidebar.bg}
\setbeamercolor{subsection in sidebar}{parent=section in sidebar}
\setbeamercolor{subsubsection in sidebar}{parent=subsection in sidebar}
\setbeamercolor{item}{fg=GSHSblue}
\setbeamercolor{block title}{bg=GSHSBLUE,fg=white}
\setbeamercolor{block body}{bg=gshsblue,fg=black}

%% Frametitle Setting %%
\setbeamerfont{frametitle}{series=\bfseries}
\setbeamertemplate{frametitle}
{ \vspace*{-10mm}
  \leavevmode
  \hspace*{3pt}
  \begin{beamercolorbox}[wd=\paperwidth,ht=1ex,dp=1ex]{frametitle}
    \hspace*{7pt}\underline{\makebox[0.6\paperwidth][l]{
    \Large{\insertframetitle}}}
  \end{beamercolorbox}
}

             
%% Title Page Setting %%
\setbeamerfont{title}{series=\bfseries}
\setbeamertemplate{title page}{
  %% Background Logo
  \begin{picture}(0,0)%
    \setlength{\unitlength}{1cm}% default
    \protect\put(1.3,0.8){%
    \begin{picture}(6,6)(4,10)%
    %%%%%%%%%%%%%%%%%%%%%%%%%%%%%%%%%%%%%%%%%%%%%
      %\includegraphics[width=0.4\paperwidth]{New_Profile_Picture.pdf}%
      \includegraphics[width=0.7\paperwidth]{koistudy_icon.pdf}%
    \end{picture}%
    }%
  \end{picture}%
  \vfill
  \vspace*{10mm}
  \raggedleft
  %% Title
  \usebeamerfont{title}\usebeamercolor[fg]{title}\inserttitle\par
  \vskip 2mm
  %% Subtitle
  \ifx\insertsubtitle\@empty
  \else\usebeamerfont{subtitle}\usebeamercolor[fg]{subtitle}\insertsubtitle
  \fi

  \vskip 3mm
  %% Horizontal line
  \usebeamercolor[fg]{title}\hrule height 2pt\hfill
  \vskip 10mm
  %% Author
  \usebeamercolor[fg]{author}\usebeamerfont{author}\insertauthor
  \vskip 1cm
  %% Institute
  \usebeamercolor[fg]{institute}\usebeamerfont{institute}\insertinstitute
  \vskip 1cm
  %% Date
  \usebeamercolor[fg]{date}\usebeamerfont{date}\insertdate
  \vfill
}

%% itemize bullet setting %%
\setbeamertemplate{itemize items}[ball]

%% block setting %%
\setbeamertemplate{blocks}[rounded][shadow=true]

%% caption setting %%
\setbeamertemplate{caption}[numbered]

%% Title, Author, Institute, Date %%
\title[O주차 - X, X]{O주차 - X, X}
\subtitle[]{2016 이산구조}
\author[]{gs1XXXX(XX기 XXX)}
\institute[GSHS]{경기과학고등학교}
\date[]{\today}

%%%%%%%%%%%%%%%%%%%%%%%%%%%%%%%%%%%%%%%%%
% insert your own profile picture
%\logo{\includegraphics[width=15mm]{stackoverflow.pdf}}
%\logo{\includegraphics[width=8mm]{Profile_Picture.pdf}}

%% Main %%

\begin{document}

\begin{frame}[plain]
\titlepage
\end{frame}
\setbeamertemplate{sidebar left}[sidebar theme]

\setstretch{1.3} %줄간격


\section{1}
\begin{frame}{1.1}\footnotesize
	관성계 : 절대 좌표계 $ (x, y, z)$ \\
	비관성계 : 회전좌표계 $ (x^{\prime}, y^{\prime}, z^{\prime})$ ~~ 극좌표 $ (r, \theta, z)$\\
	\\
	$ (x, y, z) 	\rightarrow (r, \theta, z)$\\
	수평 방향은 정역학 평형 상태에 있으므로, 	$ (x, y) 	\rightarrow (r, \theta)$\\
	$ \overrightarrow {F} = F_{x} \hat{i}  + F_{y} \hat{j} $\\
	$ F_{x} = m \frac{d^{2}x}{dt^{2}}$, $ F_{y} = m \frac{d^{2}y}{dt^{2}}$\\
	\\
	$ (x, y) = (r \cos \theta, r \sin \theta)$\\
	$ F_{r} = F_{x} \cos \theta + F_{y} \sin \theta = m \left ( \frac{d^{2}x}{dt^{2}} \cos \theta + \frac{d^{2}y}{dt^{2}} \sin \theta \right) $\\
	$ F_{\theta} = F_{y} \cos \theta - F_{x} \cos \theta = m \left ( \frac{d^{2}y}{dt^{2}} \cos \theta - \frac{d^{2}x}{dt^{2}} \cos \theta \right) $\\
\end{frame}

\begin{frame}{1.2}\footnotesize
	$ x = r \cos \theta $, $\frac{dx}{dt} = \cos \theta \frac{dr}{dt} - r \sin \theta \frac{d\theta}{dt}$, \\
	$\frac{d^{2}x}{dt^{2}} = \cos \theta \frac{d^{2}r}{dt^{2}} - \sin \theta \frac{dr}{dt} - \sin \theta \frac{dr}{dt} \frac{d\theta}{dt} -r \cos \theta \frac{d^{2}\theta}{dt^{2}}$\\
	\\
	$ y = r \sin \theta $, $\frac{dy}{dt} = \sin \theta \frac{dr}{dt} + r \cos \theta \frac{d\theta}{dt}$, \\
	$\frac{d^{2}t}{dt^{2}} = \sin \theta \frac{d^{2}r}{dt^{2}} + \cos \theta \frac{dr}{dt} + \cos \theta \frac{dr}{dt} \frac{d\theta}{dt} -r \sin \theta \frac{d^{2}\theta}{dt^{2}}$\\
	\\

	정리하면, \\
	
	$ F_{r} = m \left[ \frac{d^{2}r}{dt^{2}} - r \left( {\frac{d \theta}{dt}} \right)^{2} \right] $\\
	\\
	$ F_{\theta} = m \left[ r \frac{d^{2}\theta}{dt^{2}} + 2 \frac{dr}{dt} \frac{d\theta}{dt}  \right] $\\

\end{frame}


\begin{frame}{1.3}\footnotesize

	$\frac{d \theta}{dt} = \Omega $, $ r\frac{d \theta}{dt} = r \Omega $, $u_{\theta} $, $\frac{dr}{dt} = v_{r}$\\

	$\frac{d}{dt} r^{2}\Omega = 2 r \frac{dr}{dt} \Omega + r^{2} \frac{d\Omega}{dt}\\
	 =  r \left( r \frac{d\Omega}{dt} +2 \frac{dr}{dt} \Omega \right) \\
\\
$ y = r \sin \theta $, $\frac{dy}{dt} = \sin \theta \frac{dr}{dt} + r \cos \theta \frac{d\theta}{dt}$, \\
$\frac{d^{2}t}{dt^{2}} = \sin \theta \frac{d^{2}r}{dt^{2}} + \cos \theta \frac{dr}{dt} + \cos \theta \frac{dr}{dt} \frac{d\theta}{dt} -r \sin \theta \frac{d^{2}\theta}{dt^{2}}$\\
\\

정리하면, \\

$ F_{r} = m \left[ \frac{d^{2}r}{dt^{2}} - r \left( {\frac{d \theta}{dt}} \right)^{2} \right] $\\
\\
$ F_{\theta} = m \left[ r \frac{d^{2}\theta}{dt^{2}} + 2 \frac{dr}{dt} \frac{d\theta}{dt}  \right] $\\

\end{frame}



\begin{frame}{1.4}\footnotesize

$\frac{d \theta}{dt} = \Omega $, $ r\frac{d \theta}{dt} = r \Omega $, $u_{\theta} $, $\frac{dr}{dt} = v_{r}$\\

$\frac{d}{dt} r^{2}\Omega = 2 r \frac{dr}{dt} \Omega + r^{2} \frac{d\Omega}{dt}\\
=  r \left( r \frac{d\Omega}{dt} +2 \frac{dr}{dt} \Omega \right) \\
\\
$ y = r \sin \theta $, $\frac{dy}{dt} = \sin \theta \frac{dr}{dt} + r \cos \theta \frac{d\theta}{dt}$, \\
$\frac{d^{2}t}{dt^{2}} = \sin \theta \frac{d^{2}r}{dt^{2}} + \cos \theta \frac{dr}{dt} + \cos \theta \frac{dr}{dt} \frac{d\theta}{dt} -r \sin \theta \frac{d^{2}\theta}{dt^{2}}$\\
\\

정리하면, \\

$ F_{r} = m \left[ \frac{d^{2}r}{dt^{2}} - r \left( {\frac{d \theta}{dt}} \right)^{2} \right] $\\
\\
$ F_{\theta} = m \left[ r \frac{d^{2}\theta}{dt^{2}} + 2 \frac{dr}{dt} \frac{d\theta}{dt}  \right] $\\

\end{frame}



\begin{frame}{4}
	
	\begin{block}{asdf}
		Lorem Ipsum Dolor Sit Amet
	\end{block}
\end{frame}
>>>>>>> 7f0e4e0c82aa00f80d9da7c2f48186996b2060ff
\end{document}